\documentclass[10pt,a4paper]{amsart}
\usepackage[margin=19mm]{geometry}
\usepackage{amsmath,amssymb,mathtools}
\usepackage[scr=rsfs,cal=euler]{mathalfa}
\usepackage{xfrac}
\usepackage[unicode,hidelinks,bookmarks=false]{hyperref}
\newcommand{\ie}{i.e.\ }
\newcommand{\cf}{cf.\ }
\newcommand{\resp}{respectively\ }
\newcommand{\Subsection}{\subsection}
\providecommand{\nobreakdash}{\nobreak\hbox{-}}
\setlength{\emergencystretch}{2em}
\multlinegap0pt
\usepackage{bm}
\usepackage{bbm}
\usepackage{dsfont}
\usepackage{xspace}
\usepackage{cancel}
\usepackage{mathtools}
\usepackage{amsthm}
\makeatletter
\@ifundefined{prop}{\newtheorem{prop}{Proposition}}{}
\@ifundefined{thm}{\newtheorem{thm}{Theorem}}{}
\@ifundefined{thma}{\newtheorem{thma}{Theorem}}{}
\makeatother
\makeatletter
\expandafter\ifx\csname dedication\endcsname\relax
\newenvironment{dedication}
{\vspace{0ex}\begin{quotation}\begin{center}\begin{em}}
        {\par\end{em}\end{center}\end{quotation}}
\fi
\makeatother
\title[A note on rank zero quadratic twists of a Mordell curve]{A note on rank zero quadratic twists of a Mordell curve}
\author{Ankurjyoti Chutia, Azizul Hoque, Jyotishman Kalita}
\date{Source: arXiv:1904.02073v3 (CC BY 4.0)}
\begin{document}
\maketitle

\noindent\emph{Source and licence.} A note on rank zero quadratic twists of a Mordell curve. Version arXiv:1904.02073v3 (CC BY 4.0), distributed under the Creative Commons Attribution 4.0 International licence, as recorded in the arXiv licence metadata for this exact version and shown on the abstract page https://arxiv.org/abs/1904.02073v3. Full rights notice: licenses/arxiv-1904.02073-CC-BY-4.0.txt.\par\smallskip

\noindent\textit{Editorial scope.} Counted unit: the Prop whose source label is \texttt{prop1} in the article (source0.tex lines 94--96, source label \texttt{prop1}), delivered together with the article's own complete original proof of it (source0.tex lines 98--192, one \texttt{proof} environment of 95 lines). No mathematics is written, repaired or re-derived: statement and proof are the article's own text line for line. Required context retained uncounted: thm (lines 70--72); thmsc (lines 76--77); propcfu (lines 83--85). Every definition, display and label of those lines is retained unchanged. Omitted from this package and not redistributed: 1--69, 73--75, 78--82, 86--93, 97--97, 193--559. Nothing inside a retained excerpt is omitted or altered: every retained body is delivered exactly as the article's lines are written, so no omission receipt of any kind accompanies this package. Citations: the retained excerpts carry no citation command at all, so no bibliography entry of the article is transcribed and no reference is redistributed. Reference adaptation: none was needed and none was made; every reference inside the delivered unit resolves inside it, so no unresolved reference or citation is produced. Printed slips kept literal: no printed word, symbol, hypothesis or number of the counted statement or of its proof was altered. Layout and class: the article's own class is \texttt{amsart} with packages including amsmath, amsfonts, amssymb, mathrsfs, longtable, hyperref, graphicx; the wrapper is the lane's pinned \texttt{amsart} editorial wrapper at 10pt on A4, which restates the theorem-like environments the retained text needs and adds the article's own macro definitions that the retained text needs (none), with extra wrapper packages (bm, bbm, dsfont, xspace, cancel, mathtools). The delivered numbering is the unit's own. Assets: no retained span contains a figure environment, a graphics-inclusion command, a PDF-page inclusion, a table or a TikZ picture; no image file is copied into this package, the package declares no asset dependency and no image file is redistributed with it. Licence: version arXiv:1904.02073v3 of this article is distributed under the Creative Commons Attribution 4.0 International licence, as recorded in the arXiv licence metadata for this exact version and shown on the abstract page https://arxiv.org/abs/1904.02073v3. A full rights notice is delivered at licenses/arxiv-1904.02073-CC-BY-4.0.txt. Characters: every retained excerpt is pure ASCII, so the delivered engine, XeLaTeX with its default Latin Modern text font, reports no missing character.
\begin{thm}\label{thm}
Let $D\equiv 2\pmod 3$ be a square-free odd positive integer and $3$ does not divide the class number of $\mathbb{Q}(\sqrt{-2D})$, then the ranks of both $E_{-D}:~ y^2=x^3-2D^3$ and $E_{3D}:~ y^2=x^3+54D^3$ are zero. Assuming Cohen-Lenstra heuristics, there are infinitely many such $D$ for which the ranks of $E_{-D}$ and $E_{3D}$ are zero.
\end{thm}   

\begin{thma}\label{thmsc} Let $D>1$ be a square-free integer. Let $r$ and $s$ be the $3$-ranks of the class groups of the imaginary quadratic field $\mathbb{Q}(\sqrt{-D})$ and the real quadratic field $\mathbb{Q}(\sqrt{3D})$. Then $s\leq r\leq s+1$.
\end{thma}

 \begin{prop}\label{propcfu}
 Let $D$ be as in Theorem \ref{thm}. Then $3$ does not divide the coefficients of the fundamental unit of $\mathbb{Q}(\sqrt{6D})$. 
\end{prop}

\begin{prop}\label{prop1}
Let $D$ be as in Theorem \ref{thm}. Then $$\#\{(x,y)\in E_{-D}(\mathbb{Q}): \text{ord}_p(y)\leq 0,~\text{ for all }  p\mid 6D,~ p\text{ prime}\}=0.$$
\end{prop}

\begin{proof}
To prove this proposition, it is sufficient to show that the equation 
\begin{equation}\label{eq2.1}
y^2=x^3-2D^3z^6
\end{equation}
has no integer solutions in $x,y,z$ with $\gcd(x, y, z)=1$, $\gcd(y,D)=1$ and $z\ne 0$. Without loss of generality, we assume that $(x,y,z)$ is an integer solution of \eqref{eq2.1} such that $y$ and $z$ are positive as well as $z$ is minimal. We can exclude the cases where one or both of $x$ and $y$ are even since these cases would imply that $z$ is even too, which is a contradiction. Thus the only remaining possibility is that both $x$ and $y$ are odd.

Since $D$ is square-free and $\gcd(y,D)=1$, so that $\gcd(x,2D)=1$ and $\gcd(y,2D)=1$. 
We now rewrite \eqref{eq2.1} as follows:
\begin{equation}\label{eq2.2}
\left(y+Dz^3\sqrt{-2D}\right)\left(y-Dz^3\sqrt{-2D}\right)=x^3.
\end{equation} 
Utilizing $\gcd(x,y,z)=1$ and $\gcd(x,2D)=1$, we observe that $\gcd(y+Dz^3\sqrt{-2D}, y-Dz^3\sqrt{-2D})=1$. 
Therefore \eqref{eq2.2} gives 
\begin{equation}\label{eq2.3}
y+Dz^3\sqrt{-2D}=\left(a+b\sqrt{-2D}\right)^3
\end{equation}
for some integers $a$ and $b$ satisfying $\gcd(a,b)=1$ and 
\begin{equation}\label{eq2.4}
a^2+2b^2D=x.
\end{equation}
This shows that $a$ is odd as $x$ is odd. 
Equating the real and imaginary part from \eqref{eq2.3}, we get
\begin{equation}\label{eq2.5}
y=a^3-6ab^2D,
\end{equation}
\begin{equation}\label{eq2.6}
Dz^3=3a^2b-2b^3D.
\end{equation}
Since $D\ne 3$ and it is square-free, so that reading \eqref{eq2.6} modulo $D$, we get $ab\equiv 0\pmod D$. Now \eqref{eq2.4} together with the fact that $\gcd(x,D)=1$ implies $\gcd(a, D)$=1. Thus $b\equiv 0\pmod D$ and we write $b=Db_1$ for some integer $b_1$. Hence \eqref{eq2.6} gives
\begin{equation}\label{eq2.7}
z^3=b_1(3a^2-2b_1^2D^3).
\end{equation}
Reading \eqref{eq2.7} modulo $3$, we get $z\equiv b_1D\pmod 3$.
Utilizing this in \eqref{eq2.7} and then reading modulo $9$, we get
$$b_1^3D^3\equiv 3a^2b_1-2b_1^3D^3\pmod 9.$$
This implies $b_1(a^2-b_1^2D^3)\equiv 0\pmod 3$. If $a^2-b_1^2D^3\equiv 0\pmod 3$, then $1-2\equiv 0 \pmod 3$, since $D\equiv 2\pmod 3$.  Therefore $3\mid b_1$ and thus we write $b_1=3b_2$ for some integer $b_2$. We utilize this in \eqref{eq2.7} to get $3\mid z$, and put $z=3z_1$. Therefore \eqref{eq2.7} takes the form:
\begin{equation}\label{eq2.8}
3z_1^3=b_2(a^2-6b_2^2D^3).
\end{equation}
Since $\gcd(a, b_2)=\gcd(a,b)=1$ and $3\nmid a$, so that \eqref{eq2.8} gives $3\mid b_2$. We put $b_2=3b_3$ for some integer $b_3$. Thus \eqref{eq2.8} becomes
\begin{equation*}%\label{eq2.9}
z_1^3=b_3(a^2-54b_3^2D^3).
\end{equation*}
It is clear that $\gcd(b_3, a^2-54b_3^2D^3)=1$ since $\gcd(a,b_3)=1$ due to $\gcd(a,b)=1$. Thus there exist two integers $A$ and $B$ such that $b_3=B^3$ and $a^2-54b_3^2D^3=A^3$. These together give rise to 
\begin{equation}\label{eq2.9}
a^2-54B^6D^3=A^3.
\end{equation} 
It is clear that $3\nmid A$; otherwise $3\mid a$ which is a contradiction as $3\mid b$ too but $\gcd(a,b)=1$. Since $a$ is odd, so that $A$ is odd too. Also if a prime $p\mid \gcd(a,A)$ then by \eqref{eq2.9}, $p\mid D$ and thus by \eqref{eq2.4} we get $p\mid x$ which contradicts to $\gcd(x, D)=1$. Therefore $\gcd(a, A)=1$.

We now rewrite \eqref{eq2.9} as follows:
$$\left(a+3DB^3\sqrt{6D}\right)\left(a-3DB^3\sqrt{6D}\right)=A^3.$$
It is clear that $\gcd\left(a+3DB^3\sqrt{6D},a-3DB^3\sqrt{6D}\right)=1$ since $\gcd(a,2A)=1$.
Since the class number of $\mathbb{Q}(\sqrt{-2D})$ is not divisible by $3$, so that by 
Theorem \ref{thmsc} the class number of $\mathbb{Q}(\sqrt{6D})$ is not divisible by $3$. Therefore, 
\begin{equation}\label{eq2.10}
a+3DB^3\sqrt{6D}=u\left(\alpha+\beta\sqrt{6D}\right)^3,
\end{equation}
where $u$ is unit in $\mathbb{Q}(\sqrt{6D})$ and $\alpha, \beta$ are integers such that $\gcd(\alpha, \beta)=1$ as $\gcd(a, 3BD)=1$. Since $6D\equiv 2\pmod 4$, so that the fundamental unit, $\varepsilon$ in $\mathbb{Q}(\sqrt{6D})$ is of the form $ \varepsilon=T+U\sqrt{6D}$. Therefore the only possibilities for $u$ are $1, \varepsilon$ and $\varepsilon^2$ since the higher powers of $\varepsilon$ can be  absorbed in $(\alpha+\beta\sqrt{6D})^3$. 

We first assume that $u=1$. Then \eqref{eq2.10} gives:
\begin{equation}\label{eq2.11}
a=\alpha^3+18\alpha\beta^2D,
\end{equation} 
\begin{equation}\label{eq2.12}
DB^3=\alpha^2\beta+2\beta^3 D.
\end{equation}
Since $D$ is square-free, so that \eqref{eq2.12} shows that $D\mid \alpha\beta$. Clearly, $\gcd(D,\alpha)=1$; otherwise by \eqref{eq2.11} $\gcd(D,\alpha)\mid a$ which contradicts to $\gcd(a,D)=1$. Thus $D\mid \beta$, and we put $\beta=D\beta_1$ for some rational integer $\beta_1$. Therefore \eqref{eq2.12} becomes 
$$B^3=\beta_1(\alpha^2+2\beta_1^2D^3).$$
It is clear that $\gcd(\beta_1, \alpha^2+2\beta_1^2D^3)=1$ as $\gcd(\alpha, \beta_1)=1$. Therefore there exist two integers $\beta_2$ and $\alpha_1$ such that $\beta_1=\beta_2^3$ and $\alpha^2+2\beta_1^2D^3=\alpha_1^3$. These further imply 
$$\alpha^2=\alpha_1^3-2D^3\beta_2^6.$$
This shows that $(\alpha_1, \alpha, \beta_2)$ is another solution of $\eqref{eq2.1}$ with $\gcd(\alpha,D)=1$ and $\beta_2\ne 0$. Furthermore, 
$$|\beta_2|=|\beta_1|^{\frac{1}{3}}<|B|=|b_3|^{\frac{1}{3}}=|\frac{b_2}{3}|^{\frac{1}{3}}=|\frac{b_1}{9}|^{\frac{1}{3}}<b_1|^{\frac{1}{3}}.$$
Also, \eqref{eq2.7} gives us  
$|z|=|b_1(3a^2-2b_1^2D^3)|^{\frac{1}{3}}>|b_1|^{\frac{1}{3}}$, and thus $|\beta_2|<|z|$. This contradicts the minimality of $z$. 

We now consider the case where $u=\varepsilon$. Then \eqref{eq2.10} gives the following:
\begin{equation}\label{eq2.13}
a=T\alpha(\alpha^2+18\beta^2D)+18DU\beta(\alpha^2+2D\beta^2),
\end{equation} 
\begin{equation}\label{eq2.14}
3DB^3=3T\beta(\alpha^2+6\beta^2D)+U\alpha(\alpha^2+18D\beta^2).
\end{equation}
We read \eqref{eq2.13} modulo $3$ to get $\alpha^3T\equiv a\pmod 3$. This imples $3\nmid\alpha$  and $3\nmid T$ as $3\nmid a$. Reading \eqref{eq2.14} modulo $3$, we get 
$U\alpha^3\equiv 0\pmod 3$. This implies $3\mid U$ since $3\nmid \alpha$. This contradicts to Proposition \ref{propcfu}.

Finally if $u=\varepsilon^2$, then \eqref{eq2.10} gives 
\begin{equation}\label{eq2.15}
a=(T^2+6DU^2)(\alpha^3+18\alpha\beta^2D)+12DTU(3\alpha^2\beta+6D\beta^3),
\end{equation}
\begin{equation}\label{eq2.16}
3DB^3=(T^2+6DU^2)(3\alpha^2\beta+6D\beta^3)+2TU(\alpha^3+18\alpha\beta^2D).
\end{equation}
Reading \eqref{eq2.15} modulo $3$, we get $3\nmid T$ and $3\nmid \alpha$ since $3\nmid a$. We finally read \eqref{eq2.16} to get $TU\alpha\equiv 0\pmod 3$ which implies $U\equiv 0\pmod 3$. This again contradicts to Proposition \ref{propcfu}. 
\end{proof}

\end{document}
